\documentclass[11pt]{article}
\usepackage[a4paper,margin=1in]{geometry}
\usepackage{amsmath,amssymb,amsthm}
\usepackage{enumitem}
\usepackage[hidelinks]{hyperref}

\newtheorem{theorem}{Theorem}
\newtheorem{lemma}{Lemma}
\newcommand{\Z}{\mathbb Z}
\newcommand{\Tr}{\operatorname{Tr}}

\title{A continuous squared-phase curve through the dimension-five SIC}
\author{FA/Banach run packet for arXiv:1903.06721}
\date{13 August 2026}

\begin{document}
\maketitle

\section*{Source conjecture and result}

Appleby--Bengtsson--Flammia--Goyeneche conjecture in Section 6.5 of
arXiv:1903.06721 that the squared-phase matrix for every SIC in dimension
greater than 2 lies on a continuous curve. They construct curves in
dimensions \(3,4,6,8\), and report restricted defect \(1\) for the unique
known dimension-five orbit \(5a\), but give no dimension-five curve.

This packet proves that missing case.

\begin{theorem}\label{thm:main}
The squared-phase matrix associated with the exact Weyl--Heisenberg SIC on
Scott--Grassl orbit \(5a\) lies on a nonconstant real-analytic one-parameter
curve of squared-phase matrices satisfying equations (52)--(54) of
arXiv:1903.06721. Consequently, all six associated structures in that paper
(the projector, Hermitian Hadamard matrix, two ETFs, and two symmetric tight
fusion frames) also lie on continuous curves of structures of the same type.
\end{theorem}

This is a partial result for the source's all-SIC conjecture: it settles orbit
\(5a\), not arbitrary dimensions or hypothetical SICs outside the known
Weyl--Heisenberg orbit.

\section*{Weyl operators and the four-phase ansatz}

Work over \(\Z_5\), put \(\tau=-e^{\pi i/5}=e^{6\pi i/5}\), and let
\[
 D_{(a,b)}e_j=\tau^{b(a+2j)}e_{j+a}\qquad(a,b,j\in\Z_5).
\]
For unit complex numbers \(u,v,w,z\), define
\begin{equation}\label{eq:pattern}
 M(u,v,w,z)=
 \begin{pmatrix}
 1&u&v&\bar v&\bar u\\
 u&\bar u&w&z&w\\
 v&\bar w&\bar v&z&z\\
 \bar v&\bar z&\bar z&v&w\\
 \bar u&\bar w&\bar z&\bar w&u
 \end{pmatrix}.
\end{equation}
Thus \(M_0=1\), \(|M_p|=1\), and \(M_{-p}=\overline{M_p}\).

The exact orbit-\(5a\) fiducial in Scott--Grassl, arXiv:0910.5784,
Appendix B has unnormalized vector \(\phi\) expressed by nested square roots
(the source-archive file is \nolinkurl{sicsymbolic_5a.txt}). Set
\[
 M^\phi_p=
 \begin{cases}
 1,&p=0,\\
 6\bigl(\langle\phi,D_p^*\phi\rangle/
               \langle\phi,\phi\rangle\bigr)^2,&p\ne0.
 \end{cases}
\]
An exact finite Clifford/displacement normalization is
\begin{equation}\label{eq:normalization}
 \widetilde M_{(x,y)}
 =\tau^{-4(4y-x)}M^\phi_{(3x+y,\,2x+y)}.
\end{equation}
Direct substitution of the published radicals gives
\(\widetilde M=M(u_0,v_0,w_0,z_0)\). Numerically (only for readability),
\begin{align*}
u_0&=-0.073937094602198+0.997262907182347i,\\
v_0&= 0.877119421692956-0.480272339500221i,\\
w_0&=-0.272762046860239+0.962081527622588i,\\
z_0&=-0.860547299152712-0.509370538921297i.
\end{align*}
Equation~\eqref{eq:normalization} preserves the squared-phase equations: its
linear part has determinant one modulo \(5\), and its prefactor is a
character. The finite radical substitution and the pattern identities are
reproduced in the verifier accompanying this packet.

\section*{The \(1+2+2\) symmetry reduction}

Let \(F_0=\operatorname{diag}(1,3)\) over \(\Z_5\) and define
\begin{equation}\label{eq:A}
 A(u,v,w,z)=\frac15\sum_{p\in\Z_5^2}M(u,v,w,z)_{F_0p}D_p.
\end{equation}
This is Hermitian. Orthogonality of the Weyl operators and unimodularity of
the coefficients give, identically throughout the four-torus,
\begin{equation}\label{eq:automatic}
 \Tr A=1,\qquad \Tr(A^2)=5.
\end{equation}

The pattern~\eqref{eq:pattern} is invariant under
\(T=\left(\begin{smallmatrix}0&-1\\1&-1\end{smallmatrix}\right)\).
Consequently \(A\) commutes with the order-three unitary
\begin{equation}\label{eq:U}
 U_{jk}=\frac1{\sqrt5}\tau^{2j^2+4jk},
 \qquad
 UD_pU^*=D_{Sp},\quad
 S=F_0^{-1}TF_0=
 \begin{pmatrix}0&2\\2&4\end{pmatrix}.
\end{equation}
A direct finite geometric-sum calculation shows that the three eigenspaces of
\(U\) have dimensions \(1,2,2\). Put \(\omega=e^{2\pi i/3}\) and let
\[
 P_0=\frac{I+U+U^2}{3},\qquad
 P_1=\frac{I+\omega^2U+\omega U^2}{3};
\]
the omitted spectral projector \(P_2\) also has rank \(2\).

Write \(u=e^{ia},v=e^{ib},w=e^{ic},z=e^{id}\), and define the real-analytic map
\(f:\mathbb R^4\to\mathbb R^3\) by
\begin{equation}\label{eq:f}
 f(a,b,c,d)=
 \left(
 \Tr(P_0A)-1,\,
 \Tr(P_1A),\,
 \Tr(P_1A^2)-2
 \right).
\end{equation}

\begin{lemma}\label{lem:block}
If \(f(a,b,c,d)=0\), then \(A(a,b,c,d)^2=I_5\).
\end{lemma}

\begin{proof}
Since \(A\) commutes with \(U\), it is a direct sum of Hermitian blocks
\(A_0,A_1,A_2\) of sizes \(1,2,2\). The first equation gives \(A_0=(1)\).
The next two give \(\Tr A_1=0\) and \(\Tr(A_1^2)=2\), hence
\(A_1^2=I_2\). The automatic identities~\eqref{eq:automatic} now imply
\(\Tr A_2=0\) and \(\Tr(A_2^2)=2\), hence \(A_2^2=I_2\) as well. (For a
Hermitian \(2\times2\) matrix, trace zero and squared trace two force its
eigenvalues to be \(1,-1\).)
\end{proof}

\section*{A rigorous Jacobian certificate}

At the exact radical point
\(x_0=(\arg u_0,\arg v_0,\arg w_0,\arg z_0)\), the SIC equations and the
equivalence lemma in Section 6 of arXiv:1903.06721 give \(A(x_0)^2=I_5\), so
\(f(x_0)=0\). Differentiating~\eqref{eq:f}, an outward-rounded interval
calculation from the exact radicals gives
\[
 Df(x_0)\approx
 \begin{pmatrix}
 -0.645442112& 0.310838788& 0.237839634&-0.125923322\\
 -0.083393887&-0.351000491&-0.752846190&-0.272668334\\
  1.302407962& 0.889323925&-0.018286153&-1.338447426
 \end{pmatrix}.
\]
More importantly, the verifier certifies
\begin{equation}\label{eq:cert}
 \det D_{(a,c,d)}f(x_0)\in
 [-0.881829962934411,-0.881829962934410].
\end{equation}
All operations in this certificate are rational arithmetic, principal square
roots, and outward-rounded interval operations; the code contains the full
Scott--Grassl radical vector rather than the displayed decimals.

It follows from~\eqref{eq:cert} and the real-analytic implicit-function
theorem that, for \(b\) in an interval about \(b_0\), there are real-analytic
functions \(a(b),c(b),d(b)\) with
\(f(a(b),b,c(b),d(b))=0\). Because \(b\) itself varies, the phase
matrix~\eqref{eq:pattern} is nonconstant. Lemma~\ref{lem:block} gives
\(A^2=I_5\) along the curve. The equivalence lemma of the source paper then
converts this into the squared-phase equations and all six associated
families. This proves Theorem~\ref{thm:main}.

\section*{Verification, novelty, and limitations}

The script \texttt{code/verify\_interval.py} prints the four exact-base
enclosures, all twelve Jacobian-entry enclosures, and~\eqref{eq:cert}. It runs
with the repository's sandbox Python environment and uses \texttt{mpmath.iv}.
The exact normalization and symmetry matrices can also be checked by finite
sums over \(\Z_5\).

The run indexes and a bounded arXiv/web search on 13 August 2026 were checked
using the exact conjecture phrase and close orbit-\(5a\) variants. The search
found the source and restricted-defect discussions, but no explicit curve or
proof for orbit \(5a\). This is a likely-valid new partial result, not a
claimed full solution of the all-dimensional conjecture.

A deeper upgrade attempt considered arbitrary Zauner-symmetric SICs. Here all
symmetry blocks have size at most two. Larger blocks need higher spectral
conditions, but at an involution the differentials of \(\Tr(B^k)\) collapse to
combinations of \(\Tr(dB)\) and \(\Tr(B\,dB)\). Thus the larger-dimensional
system is singular and needs a genuinely second-order or explicit idea.

\section*{References}

\begin{itemize}[leftmargin=2em,nosep]
\item M. Appleby, I. Bengtsson, S. Flammia, and D. Goyeneche,
\emph{Tight Frames, Hadamard Matrices and Zauner's Conjecture},
arXiv:1903.06721.
\item A. J. Scott and M. Grassl, \emph{SIC-POVMs: A new computer study},
J. Math. Phys. 51 (2010), 042203; arXiv:0910.5784.
\end{itemize}

\medskip
\noindent\textbf{Human-review recommendation.}
Review as a likely valid substantial partial result, prioritizing the finite
normalization~\eqref{eq:normalization}, the \(1+2+2\) block argument, and the
interval determinant~\eqref{eq:cert}.

\end{document}
